\section{Minimum Number of Days to Make M Bouquets}
\label{sec:bs-bouquets}

\problemstatement{We are given an array $\textit{bloomDay}$ of $n$ flowers,
where $\textit{bloomDay}[i]$ is the day on which flower $i$ blooms. We need
to make $m$ bouquets, each requiring exactly $k$ \textbf{adjacent}
(consecutive in the array) flowers that have \emph{already bloomed}. Find
the \textbf{minimum day} on which it becomes possible to make all $m$
bouquets, or $-1$ if it is never possible (e.g.\ if $m \times k > n$).}

\begin{patternbox}
\emph{``Minimum day on which some condition becomes achievable''} is
binary search on the answer again, with the candidate now being a
\emph{day} rather than a speed (Section~\ref{sec:bs-koko}) or a numeric
root (Sections~\ref{sec:bs-sqrt}--\ref{sec:bs-nth-root}). The keyword
combination \emph{``minimum day/time such that a task becomes possible''}
is one of the most common phrasings of this pattern, and the underlying
monotonicity is intuitive: waiting longer can only help (more flowers have
bloomed), never hurt.
\end{patternbox}

\subsection*{Understanding the problem carefully}

For a fixed candidate day $d$, a flower $i$ is usable if
$\textit{bloomDay}[i] \le d$. Walking through the array in order, every
maximal run of consecutive usable flowers of length $L$ can be split into
$\lfloor L/k \rfloor$ bouquets (using groups of $k$ adjacent usable flowers
within that run; any leftover flowers in the run, fewer than $k$, cannot
form another bouquet). Summing $\lfloor L/k \rfloor$ over every maximal
usable run gives the total number of bouquets makeable by day $d$. Define
$\textit{bouquetsPossible}(d)$ as this total. As $d$ increases, no flower
that was previously usable becomes unusable, and some additional flowers
may become usable, so usable runs can only grow or merge -- meaning
$\textit{bouquetsPossible}(d)$ is non-decreasing in $d$.

\begin{ideabox}[Candidate Space and Predicate]
The candidate space is days in $[\min(\textit{bloomDay}),
\max(\textit{bloomDay})]$ (waiting before the earliest bloom or past the
latest bloom is never useful). The predicate is
$P(d) = (\textit{bouquetsPossible}(d) \ge m)$, and we want the smallest
$d$ where $P$ holds.
\end{ideabox}

\subsection*{Why merging consecutive usable flowers correctly counts bouquets}

The requirement that a bouquet's $k$ flowers be \emph{adjacent} in the
original array is what makes this feasibility check more than a simple
count of usable flowers divided by $k$: scattered usable flowers separated
by an unbloomed flower cannot be combined into one bouquet, even if their
total count would otherwise suffice. By scanning left to right and
resetting a running streak counter to $0$ whenever an unusable flower is
encountered (and incrementing it, and adding to the bouquet count
whenever the streak count reaches a multiple of $k$, otherwise just
incrementing), we correctly respect the adjacency constraint while still
computing $\textit{bouquetsPossible}(d)$ in a single $O(n)$ pass.

\subsection*{Algorithm}

\begin{enumerate}
  \item If $m \times k > n$: return $-1$ immediately (not enough flowers
        could ever exist).
  \item Initialize $\textit{lo} \gets \min(\textit{bloomDay})$,
        $\textit{hi} \gets \max(\textit{bloomDay})$.
  \item While $\textit{lo} \le \textit{hi}$:
  \begin{enumerate}
    \item $\textit{mid} \gets \textit{lo}+(\textit{hi}-\textit{lo})/2$.
    \item Compute $\textit{bouquetsPossible}(\textit{mid})$ via a single
          linear scan with a streak counter.
    \item If $\textit{bouquetsPossible}(\textit{mid}) \ge m$: feasible;
          try earlier: $\textit{hi} \gets \textit{mid}-1$.
    \item Else: infeasible; try later: $\textit{lo} \gets \textit{mid}+1$.
  \end{enumerate}
  \item Return \textit{lo}.
\end{enumerate}

\begin{algorithm}[H]
\caption{Minimum Day for M Bouquets}
\begin{algorithmic}[1]
\Procedure{MinDays}{\textit{bloomDay}, $m$, $k$}
  \If{$m \times k > n$} \Return $-1$ \EndIf
  \State $\textit{lo} \gets \min(\textit{bloomDay})$, $\textit{hi} \gets \max(\textit{bloomDay})$
  \While{$\textit{lo} \le \textit{hi}$}
    \State $\textit{mid} \gets \textit{lo}+(\textit{hi}-\textit{lo})/2$
    \State $\textit{bouquets} \gets \Call{BouquetsPossible}{\textit{bloomDay}, \textit{mid}, k}$
    \If{$\textit{bouquets} \ge m$}
      \State $\textit{hi} \gets \textit{mid} - 1$
    \Else
      \State $\textit{lo} \gets \textit{mid} + 1$
    \EndIf
  \EndWhile
  \State \Return \textit{lo}
\EndProcedure
\end{algorithmic}
\end{algorithm}

\subsection*{Dry run}

\dryrun{Take $\textit{bloomDay} = [7,7,7,7,12,7,7]$, $m=2$, $k=3$.}

\begin{itemize}
  \item $m \times k = 6 \le n = 7$, so an answer may exist.
  \item Try $d=12$ (the maximum bloom day, as an early sanity check): every
        flower is usable, giving one run of length $7$, so
        $\textit{bouquetsPossible}(12) = \lfloor 7/3 \rfloor = 2 \ge 2$.
        Feasible.
  \item Try $d=7$ (binary search would reach this candidate): flowers at
        indices $0,1,2,3,5,6$ are usable (bloom day $7 \le 7$), but index
        $4$ (bloom day $12$) is not. This splits the array into two runs:
        indices $[0,3]$ (length $4$) and $[5,6]$ (length $2$). Bouquets:
        $\lfloor 4/3 \rfloor + \lfloor 2/3 \rfloor = 1 + 0 = 1 < 2$.
        Infeasible.
\end{itemize}

The binary search narrows between $7$ and $12$; checking the remaining
candidates confirms that day $12$ is the earliest day at which $2$
bouquets become possible (the run of $4$ usable flowers before index $4$
cannot be split to free up a third group without waiting for the last
flower to bloom).

\subsection*{C++ implementation}

\begin{lstlisting}
#include <bits/stdc++.h>
using namespace std;

int bouquetsPossible(vector<int>& bloomDay, int day, int k) {
    int n = bloomDay.size();
    int streak = 0, bouquets = 0;

    for (int i = 0; i < n; i++) {
        if (bloomDay[i] <= day) {
            streak++;
            if (streak == k) {
                bouquets++;
                streak = 0;
            }
        } else {
            streak = 0;
        }
    }
    return bouquets;
}

int minDays(vector<int>& bloomDay, int m, int k) {
    int n = bloomDay.size();
    if ((long long)m * k > n) return -1;

    int lo = *min_element(bloomDay.begin(), bloomDay.end());
    int hi = *max_element(bloomDay.begin(), bloomDay.end());

    while (lo <= hi) {
        int mid = lo + (hi - lo) / 2;
        if (bouquetsPossible(bloomDay, mid, k) >= m) {
            hi = mid - 1;
        } else {
            lo = mid + 1;
        }
    }
    return lo;
}

int main() {
    vector<int> bloomDay = {7, 7, 7, 7, 12, 7, 7};
    cout << "Minimum day: " << minDays(bloomDay, 2, 3) << endl;
    return 0;
}
\end{lstlisting}

\begin{complexitybox}
\textbf{Time Complexity.} The candidate day range halves each iteration,
costing $O(\log(\max(\textit{bloomDay}) - \min(\textit{bloomDay})))$
iterations, and each feasibility check costs $O(n)$, giving an overall
time complexity of
\[
  O(n \log(\max(\textit{bloomDay}))).
\]
\textbf{Space Complexity.} $O(1)$ auxiliary space.
\end{complexitybox}

\begin{takeawaybox}
When a feasibility check involves an \emph{adjacency} or
\emph{contiguity} requirement (here, $k$ consecutive bloomed flowers),
the linear-scan feasibility function typically needs a running streak
counter that resets on disqualification, rather than a simple global
count. This adjacency-aware counting pattern recurs in
Section~\ref{sec:bs-painters-partition} and is worth recognizing as its
own small sub-template within the larger binary-search-on-the-answer
recipe.
\end{takeawaybox}
